\section{判别器可靠性及跨Session泛化分析}
\label{sec:ch06_discriminator_reliability}

迁移评价要求判别器首先能够稳定接受真实触觉。本节考察真实域接受能力、生成分数及跨Session校准差异。审计记录编号为P51，未执行训练或参数优化，$D_{50}$权重与类别方向保持冻结。

模型SHA256与登记一致，历史重放最大logit差为0，分类结果相同。Real fit为24个Session、4217窗，inner为8个Session、936窗，后者用于包络校准；开发评价使用另外4个Real Session、780窗及10个Sim Root、905窗。分组无交叉，但外层已有开发暴露，不能称为独立盲测。固定Real为1、Sim为0，概率阈值仍为0.5，即logit阈值0；Joint要求三支同时接受，连续Joint分数取三支logit最小值。

真实域Joint接受率先在Session内计算，再作等权宏平均，结果为0.6808333333，未达到原参考目标0.90。历史Bootstrap以4个评价Session为单位有放回抽样，每次抽取4组，对相应组接受率取等权均值；共1000次，以重采样宏均值的2.5\%和97.5\%经验分位数给出95\%区间$[0.4866666667,0.875]$，上界仍低于0.90。表\ref{tab:ch06_real_session_reliability}显示，前两个Session均为87.5\%，另外两个仅约50.67\%和46.67\%。合并780个Real窗口后，566个接受、214个拒绝，窗口合并接受率$566/780=0.7256410256$。该值按窗口加权，正式评价则按Session等权。评价仅含4个Session，不确定性估计仍受会话覆盖范围限制；组内窗口主要增加同一条件下的观测。

\begin{table}[htbp]
\centering
\caption{开发Real Session的Joint接受率}
\label{tab:ch06_real_session_reliability}
\begin{tabular}{lrr}
\toprule
Real评价Session（日期与时刻） & 窗口 & Joint接受率\\
\midrule
20260814\_150537 & 240 & 87.500\%\\
20260814\_151336 & 240 & 87.500\%\\
20260816\_172352 & 150 & 50.667\%\\
20260816\_173235 & 150 & 46.667\%\\
\bottomrule
\end{tabular}
\end{table}

相比之下，$F_{50}$的Joint FR为1，即生成样本全部被判为Real；其固定标签方向AUC仅为0.02726444845，反映分数排序严重反向，不能翻转后宣称通过。记录的混淆矩阵为$\left[\begin{smallmatrix}0&905\\214&566\end{smallmatrix}\right]$，行依次为真实类别Sim、Real，列依次为预测Sim、Real；它表示905个生成窗口均被接受，同时214个真实窗口被拒绝。矩阵按窗口计数，FR与Real宏接受率则另按来源分组。生成样本越过阈值的同时，真实样本仍可能被拒绝；AUC对连续分数排序的检查，为这一阈值现象提供了补充描述。

为观察分数是否偏离Real参照，历史审计以8个inner Session等权、组内窗口等权，在Global、Spatial、Temporal原始logit及Joint最小logit上分别取2.5\%和97.5\%加权分位数。分位数通过排序累计权重插值得到。该包络描述inner样本的经验分数范围，与原0.5分类阈值承担不同作用；Joint包络由最小logit直接构造，区别于三支包络的交集。

校准稳定性采用留一采集会话诊断（Leave-One-Session-Out，LOSO），考察包络迁移到未参与校准的会话时能否保持覆盖：每次以7个inner Session重建包络，在留下的1个Session检查覆盖，遍历8组且不选择最佳划分。

\begin{table}[htbp]
\centering
\caption{各判别分支的LOSO包络覆盖率}
\label{tab:ch06_loso_envelope_coverage}
\begin{tabular}{lrr}
\toprule
分支 & 平均覆盖率 & 最低覆盖率\\
Global & 77.083\% & 4.167\%\\
Spatial & 81.124\% & 20.417\%\\
Temporal & 94.713\% & 80.000\%\\
Joint & 85.573\% & 12.917\%\\
\bottomrule
\end{tabular}
\end{table}

\begin{table}[htbp]
\centering
\caption{Real校准分数包络与生成样本上界超出比例}
\label{tab:ch06_score_envelope_diagnostic}
\begin{tabular}{lrrr}
\toprule
分支 & logit下界 & logit上界 & $F_{50}$高于上界\\
\midrule
Global & -3.871067 & 3.667660 & 95.704\%\\
Spatial & -1.904541 & 3.850504 & 97.907\%\\
Temporal & -0.573113 & 1.827313 & 84.431\%\\
Joint & -3.871067 & 1.825418 & 84.993\%\\
\bottomrule
\end{tabular}
\end{table}

表\ref{tab:ch06_score_envelope_diagnostic}中，$F_{50}$各支均有大量分数超过Real上界，比例按Sim Root宏平均。生成分数集中在Real经验范围之外。Original的Joint包络覆盖约100\%，则显示包络覆盖与域验收的分工不同。

表\ref{tab:ch06_loso_envelope_coverage}中，Temporal平均较高，但其余分支的最低覆盖与平均相差明显。Global最低发生于20260814\_143310，Spatial最低于20260813\_165855，Joint最低于20260814\_155240；Temporal最低亦为20260814\_143310。这些差异表明，某一Session上的较好覆盖在转移到另一Session后仍可能失效。诊断范围为8个历史inner Session，未登记新的稳定性硬门槛；具体差异来源还需结合传感器、任务和物体条件进一步区分。

既定Real宏接受要求未满足，P51保留NO-GO：$D_{50}$不适合作为当前实验的最终可靠迁移评价器。该判断限于现有模型与开发数据，P49的$D_3$不在审计范围。可解释的生成接受率，需要评价器对真实触觉的稳定识别作为参照。
